\documentclass[11pt]{article}

\usepackage[a4paper,margin=1in]{geometry}
\usepackage{amsmath,amssymb,mathtools}
\usepackage{hyperref}

\hypersetup{
    colorlinks=true,
    linkcolor=blue,
    citecolor=blue,
    urlcolor=blue
}

\newcommand{\R}{\mathbb R}
\newcommand{\cP}{\mathcal P}

\title{Proof of the Benamou--Brenier Formula}
\author{}
\date{}

\begin{document}

\maketitle

\section*{Statement}

Let
\[
\mu_0,\mu_1\in \mathcal P_2(\mathbb R^d).
\]
Then
\[
W_2^2(\mu_0,\mu_1)
=
\inf_{\rho_t,v_t}
\int_0^1\int_{\mathbb R^d}|v_t(x)|^2\,d\rho_t(x)\,dt,
\]
where the infimum is taken over all narrowly continuous curves
\[
t\mapsto \rho_t\in \mathcal P_2(\mathbb R^d)
\]
and Borel velocity fields $v_t\in L^2(\rho_t)$ satisfying the continuity
equation
\[
\partial_t\rho_t+\operatorname{div}(\rho_t v_t)=0
\]
in the sense of distributions, with endpoint conditions
\[
\rho_0=\mu_0,
\qquad
\rho_1=\mu_1.
\]

\section*{Proof}

Denote the dynamic action by
\[
\mathcal A(\rho,v)
=
\int_0^1\int_{\mathbb R^d}|v_t(x)|^2\,d\rho_t(x)\,dt.
\]
We prove the equality by proving the two inequalities.

\section*{Step 1: Every admissible flow costs at least $W_2^2$}

Let $(\rho_t,v_t)$ be admissible and assume that $\mathcal A(\rho,v)<\infty$.
The continuity equation says that the evolution of $\rho_t$ can be realized by
moving particles with Eulerian velocity $v_t$.

The analytic result behind this statement is the superposition principle: there
exists a probability measure $\eta$ on absolutely continuous paths
\[
\omega:[0,1]\to\mathbb R^d
\]
such that
\[
(e_t)_\#\eta=\rho_t,
\qquad
e_t(\omega)=\omega(t),
\]
and, for $\eta$-almost every path,
\[
\dot\omega(t)=v_t(\omega(t))
\]
for almost every $t\in[0,1]$.

Define the endpoint coupling
\[
\gamma=(e_0,e_1)_\#\eta.
\]
Since $(e_0)_\#\eta=\rho_0=\mu_0$ and $(e_1)_\#\eta=\rho_1=\mu_1$, we have
\[
\gamma\in \Pi(\mu_0,\mu_1).
\]
Therefore
\[
W_2^2(\mu_0,\mu_1)
\le
\int_{\mathbb R^d\times\mathbb R^d}|x-y|^2\,d\gamma(x,y).
\]
Using the definition of $\gamma$, this becomes
\[
\int_{\mathbb R^d\times\mathbb R^d}|x-y|^2\,d\gamma(x,y)
=
\int |\omega(1)-\omega(0)|^2\,d\eta(\omega).
\]
For each absolutely continuous path, Cauchy's inequality gives
\[
|\omega(1)-\omega(0)|^2
=
\left|\int_0^1 \dot\omega(t)\,dt\right|^2
\le
\int_0^1|\dot\omega(t)|^2\,dt.
\]
Integrating with respect to $\eta$ yields
\[
W_2^2(\mu_0,\mu_1)
\le
\int\int_0^1|\dot\omega(t)|^2\,dt\,d\eta(\omega).
\]
Since $\dot\omega(t)=v_t(\omega(t))$ and $(e_t)_\#\eta=\rho_t$,
\[
\int\int_0^1|\dot\omega(t)|^2\,dt\,d\eta(\omega)
=
\int_0^1\int_{\mathbb R^d}|v_t(x)|^2\,d\rho_t(x)\,dt.
\]
Thus every admissible pair satisfies
\[
W_2^2(\mu_0,\mu_1)
\le
\mathcal A(\rho,v).
\]
Taking the infimum over all admissible pairs gives
\[
W_2^2(\mu_0,\mu_1)
\le
\inf_{\rho,v}\mathcal A(\rho,v).
\]

\section*{Step 2: An optimal transport plan realizes this cost}

Let $\gamma\in\Pi(\mu_0,\mu_1)$ be an optimal coupling for the quadratic
transport problem, so that
\[
\int_{\mathbb R^d\times\mathbb R^d}|x-y|^2\,d\gamma(x,y)
=
W_2^2(\mu_0,\mu_1).
\]
For each pair $(x,y)$, move the particle from $x$ to $y$ along the straight
line
\[
X_t(x,y)=(1-t)x+ty.
\]
Define
\[
\rho_t=(X_t)_\#\gamma.
\]
This gives $\rho_0=\mu_0$ and $\rho_1=\mu_1$.

The Lagrangian velocity of the particle labelled by $(x,y)$ is constant:
\[
\frac{d}{dt}X_t(x,y)=y-x.
\]
To express this motion in Eulerian form, disintegrate the labels $(x,y)$
conditioned on the current position $z=X_t(x,y)$ and define
\[
v_t(z)
=
\mathbb E_\gamma[y-x\mid X_t(x,y)=z].
\]
Equivalently, $v_t$ is the conditional average velocity of all particles that
occupy the point $z$ at time $t$.

For every smooth compactly supported test function $\varphi$, we compute
\[
\frac{d}{dt}\int_{\mathbb R^d}\varphi(z)\,d\rho_t(z)
=
\frac{d}{dt}\int_{\mathbb R^d\times\mathbb R^d}
\varphi(X_t(x,y))\,d\gamma(x,y).
\]
By the chain rule,
\[
\frac{d}{dt}\int \varphi(X_t(x,y))\,d\gamma(x,y)
=
\int \nabla\varphi(X_t(x,y))\cdot (y-x)\,d\gamma(x,y).
\]
By the definition of the conditional expectation $v_t$,
\[
\int \nabla\varphi(X_t(x,y))\cdot (y-x)\,d\gamma(x,y)
=
\int_{\mathbb R^d}\nabla\varphi(z)\cdot v_t(z)\,d\rho_t(z).
\]
This is precisely the weak form of
\[
\partial_t\rho_t+\operatorname{div}(\rho_t v_t)=0.
\]
Hence $(\rho_t,v_t)$ is admissible.

It remains to estimate its action. By Jensen's inequality for conditional
expectations,
\[
|v_t(z)|^2
=
\left|\mathbb E_\gamma[y-x\mid X_t=z]\right|^2
\le
\mathbb E_\gamma[|y-x|^2\mid X_t=z].
\]
Integrating with respect to $\rho_t$ gives
\[
\int_{\mathbb R^d}|v_t(z)|^2\,d\rho_t(z)
\le
\int_{\mathbb R^d\times\mathbb R^d}|y-x|^2\,d\gamma(x,y).
\]
Therefore
\[
\mathcal A(\rho,v)
\le
\int_0^1
\int_{\mathbb R^d\times\mathbb R^d}|y-x|^2\,d\gamma(x,y)\,dt
=
W_2^2(\mu_0,\mu_1).
\]
Taking the infimum over admissible pairs gives
\[
\inf_{\rho,v}\mathcal A(\rho,v)
\le
W_2^2(\mu_0,\mu_1).
\]

\section*{Conclusion}

The first step gave
\[
W_2^2(\mu_0,\mu_1)
\le
\inf_{\rho,v}\mathcal A(\rho,v),
\]
while the second step gave
\[
\inf_{\rho,v}\mathcal A(\rho,v)
\le
W_2^2(\mu_0,\mu_1).
\]
Thus
\[
W_2^2(\mu_0,\mu_1)
=
\inf_{\rho_t,v_t}
\int_0^1\int_{\mathbb R^d}|v_t(x)|^2\,d\rho_t(x)\,dt.
\]

\section*{Geometric Meaning}

The static formulation chooses an optimal coupling between initial and final
positions. The Benamou--Brenier formulation inserts time between those
positions. It says that the squared Wasserstein distance is the least kinetic
energy needed to deform $\mu_0$ into $\mu_1$ through a mass-preserving flow.

\end{document}
